\section{Einleitung}

\section*{Über das Beispielunternehmen Philetairus Immobilien GmbH}
\addcontentsline{toc}{section}{Über das Beispielunternehmen Philetairus Immobilien GmbH}  

\begin{figure}[!h]
\centering
\includegraphics[width=4.5cm, height=4.5cm]{PhiletairusLogo}
\caption{Logo der Philetairus  Immobilien GmbH, }
\end{figure}

\begin{table}[ht]
    \centering
    \begin{tabularx}{\textwidth}{|l|X|X|} 
        \hline
        Firmenname & Philetairus Immobilien GmbH\\
        \hline
        Gesellschaftsform & GmbH\\
        \hline
        Branche & Immobilienindustrie\\
        \hline
        Gründungsjahr & 1979\\
        \hline
        Standorte & Wiesbaden(Hauptsitz), \\ Limburg an der Lahn, Mainz\\
        \hline
        Produzent oder Dienstleister & Dienstleister\\
        \hline
        Anzahl Mitarbeiter & 550\\
        \hline
    \end{tabularx}
    \caption{Steckbrief Philetairus  Immobilien GmbH}
    \label{tab:firmenSteckbrief}
\end{table}

\paragraph{Hauptdienstleistungen}

Die Philetairus Immobilien GmbH ist ein Komplettdienstleister im Bereich der Immobilienwirtschaft. Das Kerngeschäft umfasst die Verwaltung von Wohnungseigentum (WEG-Verwaltung) sowie die Miet- und Sondereigentumsverwaltung. \\

Weitere Dienstleistungen beinhalten technisches, kaufmännisches und infrastrukturelles Gebäudemanagement, Vermietungsmanagement, Erstellung der Betriebskostenabrechnungen sowie Hausmeisterservices.\\

\paragraph{Firmenziele}

Die Philetairus Immobilien GmbH strebt danach, einen umfassenden Service für die zu verwaltenden Immobilien zu bieten, der sowohl Eigentümer als auch Mieter zufriedenstellt. Ein weiteres Ziel ist die Förderung einer starken Kooperation der einzelnen Unternehmensbereiche, um einen zentralen Anlaufpunkt für alle Bedürfnisse der Auftraggeber zu bieten.\\

Die hier beschriebene Firma ist rein fiktiv und wurde nur für den Zweck der Habmann Abschlussarbeit erstellt und wurde im Rahmen dieses Moduls unverändert übernommen. \\
\newpage


\subsection{Projektbeschreibung}
Kurze Einführung in das Projekt, Zielsetzung und Nutzen.

\subsubsection{Einführung}

Das Projekt PhileTipTip zielt darauf ab, ein innovatives System für die Abteilung Außenarbeiten zu entwickeln. Dieses System soll in Form einer Android-App realisiert werden und dient als Kommunikationsbrücke zwischen Mietern/Eigentümern und der Grünflächenabteilung einer Immobilienverwaltung. Zusätzlich soll ein Administrations Tool für die Desktop Anwendung entstehen, die das erfassen von neuen Datensätzen für Mieter, Gebäude und Mitarbeiterdaten vereinfacht und vereinheitlicht. Für dieses Pionierprojekt sollen beide Komponenten der Erfassung von Meldungen für die Grünflächenabteilung dienen, aber perspektivisch soll mit diesem Projekt der Grundstein für eine umfassende Digitalisierung des gesamten Unternehmens gelegt werden.

\subsubsection{Zielsetzung}

Das Hauptziel des Projekts ist es, einen effizienten und benutzerfreundlichen Weg zu schaffen, um Meldungen über Probleme im Außenbereich von Immobilien zu erfassen, zu verwalten und zu bearbeiten. Insbesondere soll das System folgende Funktionen bieten:
\begin{itemize}
  \item  Ermöglichung der Erfassung von Meldungen durch Mieter und Eigentümer
  \item  Fokus auf Probleme im Zuständigkeitsbereich der Abteilung Außenarbeiten (z.B. Schädlingsbefall, Unkrautwuchs)
  \item  Effiziente Weiterleitung der Meldungen an die zuständige Abteilung
  \item  Verwaltung und Bearbeitung der eingegangenen Meldungen durch die Abteilung Außenarbeiten
\end{itemize}

\subsubsection{Nutzen}

Die Implementierung dieses Systems bringt mehrere Vorteile mit sich:
\begin{itemize}
  \item  Verbesserte Kommunikation zwischen Bewohnern und Immobilienverwaltung
  \item  Schnellere Reaktionszeiten auf Probleme im Außenbereich
  \item  Erhöhte Effizienz in der Bearbeitung von Meldungen
  \item  Bessere Übersicht und Nachverfolgbarkeit von Außenarbeiten
  \item  Steigerung der Zufriedenheit von Mietern und Eigentümern
\end{itemize}

\subsubsection{Kernfunktionen}

Die Android-App PhileTipTip wird folgende Kernfunktionen beinhalten:
\begin{itemize}
  \item  Benutzerfreundliche Oberfläche zur Erfassung von Meldungen
  \item  Möglichkeit zum Hochladen von Fotos zur besseren Dokumentation der Probleme
  \item  Kategorisierung der Meldungen (z.B. Schädlingsbefall, Unkrautwuchs)
  \item  Automatische Weiterleitung der Meldungen an die zuständige Abteilung
  \item  Statusverfolgung der Meldungen für Benutzer
  \item  Verwaltungsinterface für die Bearbeitung und Aktualisierung von Meldungen
\end{itemize}

\subsection{Zielgruppe}

Die primären Nutzer der App sind:
\begin{itemize}
  \item  Mieter von Immobilien
  \item  Eigentümer von Immobilien
  \item  Mitarbeiter der Abteilung Außenarbeiten in der Immobilienverwaltung
\end{itemize}

\subsection{Ziele des Projekts}

Durch die Entwicklung von PhileTipTip wird ein modernes, effizientes System geschaffen, das die Kommunikation und Problemlösung im Bereich der Außenarbeiten von Immobilien signifikant verbessert.

\newpage